% answers_short.tex —— 2004 年·上海卷·理科 · 简短答案速查
% 由 tools/make_short_answers.py 自 answers.tex 的 \gksolhead/\gkanswer 提取，
% 并与 2004/2004_数学真题_上海理_答案要点.md 逐题比对后生成，勿手改（改源 MD 后重跑）。

\begin{gknotice}
\textbf{使用说明}：本节只给最终答案便于快速核对，完整推导见解析版。答案为编写者独立推导，\textbf{非官方标准答案}。
\end{gknotice}

\gkpart{填空题}
\begin{itemize}[leftmargin=2.6em]
  \item[\textbf{1.}] $3$
  \item[\textbf{2.}] $(5, 0)$
  \item[\textbf{3.}] $\{1, 2, 5\}$
  \item[\textbf{4.}] $2$
  \item[\textbf{5.}] $(-2, 0) \cup (2, 5]$
  \item[\textbf{6.}] $(5, 4)$
  \item[\textbf{7.}] $\dfrac{2\sqrt{15}}{5}$
  \item[\textbf{8.}] $(x - 2)^2 + (y + 3)^2 = 5$
  \item[\textbf{9.}] $\dfrac{4}{11}$
  \item[\textbf{10.}] $a > 0,\; b \leqslant 0$
  \item[\textbf{11.}] 用代数方法研究几何问题（或：坐标法）
  \item[\textbf{12.}] ①④
\end{itemize}

\gkpart{选择题}
\begin{center}\begin{tabular}{ccccc}
\toprule
题号 & 13 & 14 & 15 & 16 \\
\midrule
答案 & B & C & A & B \\
\bottomrule
\end{tabular}\end{center}

\gkpart{解答题}
\begin{itemize}[leftmargin=2.6em]
  \item[\textbf{17.}] $(1, 7)$（或 $1 < a < 7$）
  \item[\textbf{18.}] $x = 8 - 4\sqrt{2} \approx 2.343\text{ m},\; y = 2\sqrt{2} \approx 2.828\text{ m}$
  \item[\textbf{19.}] (1) $(-\infty, -1) \cup [1, +\infty)$；(2) U/PENDING
  \item[\textbf{20.}] (1) $f(x) = x^2 + \dfrac{8}{x}$；(2) 证明见解析
  \item[\textbf{21.}] (1) 见解析版；(2) $\arccos\dfrac{\sqrt{6}}{3}$（或 $\arcsin\dfrac{\sqrt{3}}{3}$、$\arctan\dfrac{\sqrt{2}}{2}$）；(3) 存在，构造各棱长均为 $\dfrac{1}{2}$、底面相邻两边夹角为 $\alpha = \arcsin(8V)$ 的直平行六面体
  \item[\textbf{22.}] (1) U/PENDING；(2) $\dfrac{n(a^2 + b^2)}{2}$；(3) 选抛物线 $C: y^2 = 2px$（$p > 0$），取 $P_1(0, 0)$ 时，充要条件为 $d > 0$
\end{itemize}
